\documentclass[11pt,a4paper]{article}
\usepackage[utf8]{inputenc}\usepackage[T1]{fontenc}
\usepackage[provide=*,english]{babel}
\usepackage{lmodern,microtype,amsmath,amssymb,booktabs,geometry,xcolor,fancyhdr}
\usepackage[hidelinks]{hyperref}
\geometry{margin=2.3cm,headheight=14pt}
\setlength{\parindent}{0pt}\setlength{\parskip}{0.4em}
\definecolor{gcblue}{RGB}{34,73,110}\definecolor{gcgray}{RGB}{90,96,102}
\pagestyle{fancy}\fancyhf{}
\fancyhead[L]{\small\color{gcgray}GC2D: SDIRK4 S54b}
\fancyhead[R]{\small\color{gcgray}Model theory}\fancyfoot[C]{\thepage}
\setlength{\emergencystretch}{3em}
\begin{document}
\begin{center}
 {\LARGE\bfseries\color{gcblue}Five-Stage Fourth-Order SDIRK}\par
 \vspace{0.35em}{\large Skvortsov S54b non-geometric reference method}
\end{center}

\section{Butcher tableau}

For $\dot z=f(t,z)$, the S54b tableau is
\[
\begin{array}{c|ccccc}
\tfrac14&\tfrac14&0&0&0&0\\
0&-\tfrac14&\tfrac14&0&0&0\\
\tfrac12&\tfrac18&\tfrac18&\tfrac14&0&0\\
1&-\tfrac32&\tfrac34&\tfrac32&\tfrac14&0\\
1&0&\tfrac16&\tfrac23&-\tfrac1{12}&\tfrac14\\ \hline
&0&\tfrac16&\tfrac23&-\tfrac1{12}&\tfrac14
\end{array}.
\]
All five diagonal entries equal $\gamma=1/4$, so the method is singly
diagonally implicit.  The final row of $A$ equals $b^T$, making the method
stiffly accurate.

The stages are solved sequentially:
\begin{align}
 Z_i&=z_n+h\sum_{j<i}a_{ij}f(t_n+c_jh,Z_j)
      +h\gamma f(t_n+c_ih,Z_i),\\
 z_{n+1}&=z_n+h\sum_{i=1}^{5}b_i f(t_n+c_ih,Z_i).
\end{align}
Each stage therefore requires a nonlinear solve in the physical state
dimension, rather than one simultaneous system containing all stages.

\section{Fourth-order conditions}

Writing $C=\operatorname{diag}(c)$ and using componentwise powers, the
coefficients satisfy
\[
\begin{array}{llll}
b^T\mathbf 1=1,&b^Tc=\tfrac12,&b^Tc^2=\tfrac13,&b^TAc=\tfrac16,\\
b^Tc^3=\tfrac14,&b^TCAc=\tfrac18,&b^TAc^2=\tfrac1{12},&b^TA^2c=\tfrac1{24}.
\end{array}
\]
These are the Runge--Kutta conditions through order four.  For sufficiently
smooth problems the local truncation error is $\mathcal O(h^5)$ and the global
error is $\mathcal O(h^4)$.

\section{Deliberately non-geometric structure}

A Runge--Kutta method is symplectic for general canonical Hamiltonian systems
only if
\[
 b_i a_{ij}+b_j a_{ji}-b_i b_j=0
 \qquad\text{for every }i,j.
\]
The S54b coefficients do not satisfy this identity.  For example, with
$i=j=5$ the left-hand side is $1/16$, not zero.  Hence S54b is not a
symplectic Runge--Kutta method.

Symmetry would require equality with the adjoint tableau, including
\[
 c_i=1-c_{6-i},\qquad b_i=b_{6-i},\qquad
 a_{ij}+a_{6-i,6-j}=b_{6-j}.
\]
Already $c_1=1/4\ne 1-c_5=0$, so S54b is not symmetric.  These failures make
the method a useful same-order control in long-time comparisons against
geometric fourth-order methods.

\section{Newton implementation}

For a known right-hand side $r_i$, one stage solves
\[
 R_i(Z_i)=Z_i-r_i-h\gamma f(t_n+c_ih,Z_i)=0.
\]
Newton corrections satisfy
\[
 \left[I-h\gamma D_zf(t_n+c_ih,Z_i^{(k)})\right]\delta_i^{(k)}
 =-R_i(Z_i^{(k)}).
\]
For guiding-centre dynamics the implementation uses independent analytic
$2\times2$ particle Jacobians.  A centered finite-difference dense Jacobian is
available for generic dynamical systems.  Per-stage and per-step iterations,
residual evaluations, residual norms, and effective tolerances are returned as
diagnostics.

\section{Reference}

The coefficient table is the S54b method presented by L.~M.~Skvortsov,
``Diagonally implicit Runge--Kutta methods for stiff problems,''
\emph{Computational Mathematics and Mathematical Physics}, 46(12),
2110--2123 (2006), DOI:
\href{https://doi.org/10.1134/S0965542506120098}{10.1134/S0965542506120098}.

\clearpage
\section{Four state formulations and passive energy tracking}
For one planar particle the supported internal representations are
\[
\begin{array}{lll}
\text{physical} & z=(x,y) & \mathbb{R}^{2},\\
\text{physical with energy} & (z,t,\kappa) & \mathbb{R}^{4},\\
\text{duplicated} & (u,v),\quad u,v\in\mathbb{R}^{2} & \mathbb{R}^{4},\\
\text{duplicated with energy} & (u,v,t,\kappa) & \mathbb{R}^{6}.
\end{array}
\]
Classical methods use the physical rows; ABBA and BM4 use the duplicated rows.
\texttt{track\_energy} selects the energy variant. For $N$ planar particles the
respective dimensions are $2N$, $4N$, $4N$, and $6N$: one time entry and
one momentum per particle. All time entries equal the integration time. Classical full-cyclotron runs instead use $4N$ or $6N$ entries.
Accepted spatial copies coincide; only one physical copy is exported.
Every coordinate block contains $N$ values in component-major order. The
formulation validates particle times and aligns output round-off; the
integrator does not interpret coordinate indices.

The normalized passive balance is
\[
\dot\kappa=-\partial_tH(t,z),\qquad \kappa(t_0)=0,\qquad
\Delta K(t)=H(t,z(t))+\kappa(t)-H(t_0,z_0).
\]
Splitting accumulators from both copies are divided by two. The physical
Hamiltonian may vary with time; $\Delta K$ measures its computed balance.
Tracking leaves physical trajectories, nonlinear stopping rules and adaptive
acceptance unchanged. Dynamics must supply the momentum derivative, including
zero for an autonomous Hamiltonian.

Hairer's constraint is solely $u-v=0$: the reduced multiplier has $2N$ entries
and ABBA's simultaneous spatial solve has $6N$ unknowns. Neither $t$ nor
$\kappa$ is projected. The former eight-component projection is retired;
\texttt{state\_extension='fully\_extended'} raises migration guidance.

Dynamics, formulation, method and integration remain separate. The shared
state classes own layouts, diagonal embedding and history extraction in
\begin{quote}\small\path{src/simulation/formulations/state.py}\end{quote}
Each integration initializes a fresh method instance. Observers receive physical
states; energy histories are in \texttt{Solution.diagnostics}.
The \texttt{extended\_time} history has shape $(N,n_{\rm samples})$. Fixed-step shadow
samples leave main states and counters unchanged. Adaptive energy quadrature
uses accepted dense output; its backend state and Jacobian remain physical.

The five converged physical stages supply the passive momentum quadrature with the tableau weights.

\end{document}
